\documentclass[11pt]{article}

\usepackage[a4paper,margin=1in]{geometry}
\usepackage{amsmath,amssymb}
\usepackage{hyperref}

\hypersetup{
    colorlinks=true,
    linkcolor=blue,
    citecolor=blue,
    urlcolor=blue
}

\title{The Gradient Represents the Directional Derivative Through the Metric}
\author{}
\date{}

\begin{document}

\maketitle

\section*{The Problem}

At a point $p$ of a manifold, the derivative of a function is a covector:
\[
df_p:T_pM\to\mathbb R.
\]
It is a linear functional. Given a tangent vector $X_p$, it returns the number
\[
df_p(X_p),
\]
the directional derivative of $f$ in the direction $X_p$.

But gradient descent wants a vector. So we ask:

\begin{quote}
Which tangent vector contains the same information as the covector $df_p$?
\end{quote}

\section*{Representation by the Metric}

If $g$ is a Riemannian metric, then every tangent vector $Y_p$ defines a
covector by
\[
X_p\mapsto g_p(Y_p,X_p).
\]
So vectors can represent linear functionals through the metric.

The gradient is defined as the unique vector
\[
\operatorname{grad}_g f(p)\in T_pM
\]
whose metric pairing with every direction $X_p$ reproduces the directional
derivative:
\[
df_p(X_p)
=
g_p(\operatorname{grad}_g f(p),X_p)
\]
for all $X_p\in T_pM$.

This is what ``represents'' means: the covector $df_p$ and the vector
$\operatorname{grad}_g f(p)$ give the same numerical answer after the vector is
paired with $X_p$ using $g_p$.

\section*{Euclidean Example}

In $\mathbb R^n$ with the Euclidean dot product,
\[
Df(x)[v]
=
\sum_i\frac{\partial f}{\partial x^i}(x)v^i.
\]
The vector
\[
\nabla f(x)
=
\left(
\frac{\partial f}{\partial x^1},
\ldots,
\frac{\partial f}{\partial x^n}
\right)
\]
satisfies
\[
Df(x)[v]
=
\nabla f(x)\cdot v.
\]
So the Euclidean gradient is the vector representing the covector $Df(x)$
through the Euclidean metric.

\section*{Non-Euclidean Metric}

If the metric has components $g_{ij}$, the representing vector changes:
\[
(\operatorname{grad}_g f)^i
=
g^{ij}\frac{\partial f}{\partial x^j}.
\]
The covector $df$ is the same differential object, but the vector representing
it depends on the metric.

\section*{Wasserstein Analogy}

In Wasserstein geometry, the role of the metric pairing is
\[
\langle u,v\rangle_\rho
=
\int u\cdot v\,d\rho.
\]
When we write
\[
D\mathcal F(\rho)[v]
=
\langle \operatorname{grad}_{W_2}\mathcal F(\rho),v\rangle_\rho,
\]
we are using exactly the same idea: the Wasserstein gradient is the velocity
field that represents the directional derivative through the Wasserstein
metric.

\end{document}
